% Imporditud: tools/import_eksperimendid.py
% Allikas: Eesti füüsikaolümpiaadi eksperimendi ülesannete
%   kogu (Rael Kalda, 2023), ülesanne Ü143, lahendus L143.
\setAuthor{EFO žürii}
\setRound{lõppvoor}
\setYear{2013}
\setNumber{G E2}
\setDifficulty{4}
\setTopic{dünaamika}
\setVahendid{lauatennisepall, kaks puitjoonlauda}
\setPunktid{12}
\setAllikas{Eksperimendi kogu Ü143, lahendus lk 107}
\setLahendusseis{toores}

\prob{Pinksipall}
Mõõtke seisuhõõrdetegur lauatennisepalli ja joonlaua vahel. Hinnake mõõtemääramatust.

\hint

\solu
Paneme palli kahe joonlaua vahele, hoiame nende ühed otsad koos ja vajutame teistest otstest tugevalt kokku. Sellisel juhul moodustab hõõrde- ja rõhumisjõu summa pinnanormaaliga nurga $\alpha$ = arctan($\mu$). Mõõdame nurga 2$\alpha$ joonlaudade vahel (mõõtes joonlaudadade pikkuse palliga kokkupuutepunktini b ja kokkupuutepunktide vahekaudus a, $\alpha$ = arcsin(a/2b)) ning selle abil $\mu$ = tan $\alpha$ = a/ 4b2 $-$ $a_2$.
\probend
